% Date : 19/03/2022
% Version : 4
% Auteur : Alexis Brès
% Relu : Maxence Miguel-Brebion
% Relecteur : Aela Fortun

% ------------------------------------------------------------ %
% ----------------------  Fiche MCA03  ----------------------- %
% Thème : Approche énergétique en mécanique
% Classes : MPSI-MP2I-PCSI-PTSI
% ------------------------------------------------------------ %



% ======================================================= % 
% ============== Gestion de la compilation ============== %
% ======================================================= % 
\ifdefined\mainIsLoaded\else
\RequirePackage{import}
\subimport{../../}{_preambule_CdE_PC}
\begin{document}
\fi
% ======================================================= % 
% ======================================================= % 






% ============================================================================ %
%                            FICHE D'ENTRAÎNEMENT                              %
% ============================================================================ %
% ======================= Métadonnées sur la fiche =========================== %
\titreFicheEntrainement		{Approche énergétique en mécanique}
\grandTheme					{Mécanique}
\numeroFiche				{MCA03}
\uniqueID					{mIELXcueXfn} % une chaîne de caractères (a-z, A-Z) aléatoire de 10 caractères.
\nombreColonnesReponses		{2} % 1, 2, 3, etc.
% ============================================================================ %

%%%%%%%%%%%%%%%%%%%%%%%%%%%%%%%%%%%%%%%%%%%%%%%%%%%%%%%%%%%%%%%%%%%%%%%%%%%%%%%%%%%%%%%%%%%%%%
\debutFicheEntrainement                                                                      %
%%%%%%%%%%%%%%%%%%%%%%%%%%%%%%%%%%%%%%%%%%%%%%%%%%%%%%%%%%%%%%%%%%%%%%%%%%%%%%%%%%%%%%%%%%%%%%




% --------------------------------------------------------------------- %
%                              Prérequis                                %
%                             (facultatif)                              %
% --------------------------------------------------------------------- %
% ================== Métadonnées sur le prérequis ===================== %
\avecPrerequis{Y} % Y ou N
% ===================================================================== %
\begin{prerequis}
Systèmes de coordonnées. Expression de forces (poids, force de rappel). \\
Travail d'une force. Théorèmes généraux (dynamique et énergétiques).	
\end{prerequis}
% --------------------------------------------------------------------- %






% •••••••••••••••••••••••••••••••••••••••••••••••••••••••••••••••••••••••••••• %
% •••••••••••••••••••••••••••••••••••••••••••••••••••••••••••••••••••••••••••• %
\sectionFicheEntrainement{\' Energies potentielles}
% •••••••••••••••••••••••••••••••••••••••••••••••••••••••••••••••••••••••••••• %
% •••••••••••••••••••••••••••••••••••••••••••••••••••••••••••••••••••••••••••• %


% ********************************************************************* %
%                           ENTRAÎNEMENT                                % 
% ********************************************************************* %
% ================ Métadonnées sur l'entraînement ===================== %

\titreEntrainementFacultatif	{La juste formule}
\hauteurLargeurCadreReponse		{8mm}{3cm}
\dureeResolutionFacultative		{1} % 1, 2, 3 ou 4
\typeEntrainement				{QCM} % Newton ou QCM ou AN ou vide (exercice "normal")
\nombreColonnesQuestions		{1} % vide, 1, 2, 3, etc.
\avecPlusieursQuestions			{N} % Y ou N
\initialisationEntrainement
% ===================================================================== %


% =============================================================================== %
\debutEntrainement
% =============================================================================== %

% ^^^^^^^^^^^^^^^^^^^^^^^^^^^^^^^^^^^^^^ %
%               Question                 %
% ^^^^^^^^^^^^^^^^^^^^^^^^^^^^^^^^^^^^^^ %

\begin{enonce}
	On considère un point matériel de masse $m$ plongé dans le champ de pesanteur $\vv{g}$. On se place dans un repère cartésien $(\ptO,\vv{e_x},\vv{e_y},\vv{e_z})$ tel que $\vv{g}=-g\vv{e_y}$, le point $\ptO$ étant pris comme origine de l'énergie potentielle. Quelle est l'expression de l'énergie potentielle de pesanteur ?
	
	\begin{listeQCM5Colonnes}
	\item $mgx$
	\item $-mgy$
	\item $mgy$
	\item $mgz$
	\end{listeQCM5Colonnes}
\end{enonce}

\reponse{\reponseC{}}

% ************************************** %

% =============================================================================== %
\finEntrainement
% =============================================================================== %



% ********************************************************************* %
%                           ENTRAÎNEMENT                                % 
% ********************************************************************* %
% ================ Métadonnées sur l'entraînement ===================== %

\titreEntrainementFacultatif	{Plusieurs expressions d'énergie potentielle de pesanteur}
\hauteurLargeurCadreReponse		{7mm}{4cm}
\dureeResolutionFacultative		{2} % 1, 2, 3 ou 4
\typeEntrainement				{Newton} % Newton ou QCM ou AN ou vide (exercice "normal")
\basiqueEtTransversal			{Y}
\nombreColonnesQuestions		{2} % vide, 1, 2, 3, etc.
\avecPlusieursQuestions			{Y} % Y ou N
\initialisationEntrainement
% ===================================================================== %

Déterminer la fonction énergie potentielle de pesanteur d'un point matériel de masse $m$ associée aux situations suivantes :
\begin{center}
\subimport{_images/}{Energie_potentielle_pesanteur_CBD_v2.tex}
\end{center}

% =============================================================================== %
\debutEntrainement
% =============================================================================== %

% ^^^^^^^^^^^^^^^^^^^^^^^^^^^^^^^^^^^^^^ %
%               Question                 %
% ^^^^^^^^^^^^^^^^^^^^^^^^^^^^^^^^^^^^^^ %

\begin{enonce}
$E_{pp}(y)=$
\end{enonce}

\reponse{$mg(\ell-y)$}

\begin{corrige}
L'axe est ici orienté vers le bas, on a donc $E_{pp}(y)=-mgy+ K_1$. On veut $E_{pp}(\ell)=0$, d'où $K_1=mg\ell$. Finalement, on a 
$E_{pp}(y)=mg(\ell-y)$.
\end{corrige}

% ************************************** %


% ^^^^^^^^^^^^^^^^^^^^^^^^^^^^^^^^^^^^^^ %
%               Question                 %
% ^^^^^^^^^^^^^^^^^^^^^^^^^^^^^^^^^^^^^^ %

\begin{enonce}
$E_{pp}(x) =$
\end{enonce}

\reponse{$mg\left(x\sin(\alpha)-H\right)$}

\begin{corrige}
On peut raisonner de deux manières :
\begin{itemize}
\item La coordonnée verticale (axe de $\vv{g}$) $z$ est liée à $x$ par $z=x\sin(\alpha)$. On a donc $E_{pp}=mgx\sin(\alpha)+K_2$. 

L'énergie potentielle étant nulle en $z=H$, on a $E_{pp}(x)=mg(x\sin(\alpha)-H)$.

\item Dans le repère $(\ptO, \vv{e_x}, \vv{e_y})$, on a $\vv{g}=-g\sin(\alpha)\vv{e_x}-g\cos\alpha\vv{e_y}$. 

On en déduit le travail élémentaire pour un déplacement selon $x$ :
\[
\delta W = -mg\sin(\alpha) \d{} x = -\d{}(mgx\sin(\alpha) +K_2) = -\d{}E_{pp}.
\]

On en déduit que $E_{pp}(x)=mgx\sin(\alpha)+K_2$. 

L'énergie potentielle devant être nulle en $S$, qui correspond à $x=\frac{H}{\sin(\alpha)}$, on a $K_2=-mgH$, d'où le résultat.
\end{itemize}
\end{corrige}

% ************************************** %


% ^^^^^^^^^^^^^^^^^^^^^^^^^^^^^^^^^^^^^^ %
%               Question                 %
% ^^^^^^^^^^^^^^^^^^^^^^^^^^^^^^^^^^^^^^ %

\begin{enonce}
$E_{pp}(\theta) =$
\end{enonce}

\reponse{$-mgR\cos(\theta)$}

\begin{corrige}
Dans la base polaire, l'accélération de la pesanteur s'écrit $\vv{g}=g\cos(\theta)\vv{e_r}-g\sin(\theta)\vv{e_\theta}$. Donc, le travail élémentaire pour un déplacement sur le cercle (selon $\vv{e_\theta}$) est 
\[
\delta W = m\vv{g}\cdot\d{}\vv{\ptO \ptM}=-mg\sin(\theta) R \d{}\theta =-\d{}(-mgR\cos(\theta)+K_3) =-\d{}E_{pp}.
\]

On a donc $E_{pp}(\theta)=-mgR\cos(\theta)+K_3$, et comme on veut $E_{pp}(\pi/2)=0$, on a $K_3=0$. 
Ainsi, on a 
\[
E_{pp}(\theta) = -mgR\cos(\theta).
\]
\end{corrige}

% ************************************** %


% ^^^^^^^^^^^^^^^^^^^^^^^^^^^^^^^^^^^^^^ %
%               Question                 %
% ^^^^^^^^^^^^^^^^^^^^^^^^^^^^^^^^^^^^^^ %

\begin{enonce}
$E_{pp}(\psi) =$
\end{enonce}

\reponse{$mgr\big(\cos(\psi)-1\big)+E_0$}

\begin{corrige}
Fixons un axe $(\ptO z)$ vertical ascendant avec $\ptO$ au centre du cercle. L'énergie potentielle de pesanteur s'écrit alors $E_{pp}=mgz+K_4$. Or, on a $z=r\cos(\psi)$ d'où $E_{pp}=mgr\cos(\psi)+K_4$. 

La convention choisie ($E_{pp}(\psi=0)=E_0$) entraîne que 
\[
mgr\cos(0)+K_4=E_0
\quad\text{d'où}\quad
K_4=E_0-mgr.
\]
Finalement, on trouve
\[
E_{pp}=mgr\big(\cos(\psi)-1\big)+E_0.
\]
\end{corrige}

% ************************************** %

% =============================================================================== %
\finEntrainement
% =============================================================================== %





% ********************************************************************* %
%                           ENTRAÎNEMENT                                % 
% ********************************************************************* %
% ================ Métadonnées sur l'entraînement ===================== %

\titreEntrainementFacultatif	{La juste formule... le retour}
\hauteurLargeurCadreReponse		{8mm}{3cm}
\dureeResolutionFacultative		{1} % 1, 2, 3 ou 4
\typeEntrainement				{QCM} % Newton ou QCM ou AN ou vide (exercice "normal")
\nombreColonnesQuestions		{1} % vide, 1, 2, 3, etc.
\avecPlusieursQuestions			{N} % Y ou N
\initialisationEntrainement
% ===================================================================== %


% =============================================================================== %
\debutEntrainement
% =============================================================================== %

% ^^^^^^^^^^^^^^^^^^^^^^^^^^^^^^^^^^^^^^ %
%               Question                 %
% ^^^^^^^^^^^^^^^^^^^^^^^^^^^^^^^^^^^^^^ %

\begin{enonce}
	On considère un point matériel \ptM de masse $m$ astreint à se déplacer selon un axe $(\ptO y)$ horizontal. Il est attaché à un ressort de raideur $k$ et de longueur à vide $\ell_0$. L'autre extrémité du ressort est fixée en $\ptO$. 
	
	Quelle est l'expression de l'énergie potentielle élastique du point \ptM pour que celle-ci soit nulle lorsque l'allongement du ressort est nul ?
	
	\begin{listeQCM4Colonnes}
	\item $\dfrac{1}{2}ky^2$
	\item $\dfrac{1}{2} k(y-\ell_0)^2$
	\item $\dfrac{1}{2}k(y^2-{\ell_0}^2)$
	\item $-\dfrac{1}{2}k(\ell_0-y)^2$
	\end{listeQCM4Colonnes}
		\bigskip\smallskip
\end{enonce}

\reponse{\reponseB}

% ************************************** %


% =============================================================================== %
\finEntrainement
% =============================================================================== %





% ********************************************************************* %
%                           ENTRAÎNEMENT                                % 
% ********************************************************************* %
% ================ Métadonnées sur l'entraînement ===================== %

\titreEntrainementFacultatif	{Expression de l'énergie potentielle élastique}
\hauteurLargeurCadreReponse		{8mm}{6cm}
\dureeResolutionFacultative		{3} % 1, 2, 3 ou 4
\typeEntrainement				{Newton} % Newton ou QCM ou AN ou vide (exercice "normal")
\basiqueEtTransversal			{Y}
\nombreColonnesQuestions		{1} % vide, 1, 2, 3, etc.
\avecPlusieursQuestions			{Y} % Y ou N
\initialisationEntrainement
% ===================================================================== %

Déterminer la fonction énergie potentielle élastique associée aux situations suivantes, où tous les ressorts sont de longueur à vide $\ell_0$ et de raideur $k$ :

\begin{center}
\subimport{_images/}{Energie_potentielle_elastique_GBU_v4.tex}
\end{center}

% =============================================================================== %
\debutEntrainement
% =============================================================================== %

% ^^^^^^^^^^^^^^^^^^^^^^^^^^^^^^^^^^^^^^ %
%               Question                 %
% ^^^^^^^^^^^^^^^^^^^^^^^^^^^^^^^^^^^^^^ %

\begin{enonce}
$E_{pe}(y) = $
\end{enonce}

\reponse{$\frac{1}{2}k(y-\ell_0)^2-\frac{k{\ell_0}^2}{2}$}

\begin{corrige}
L'axe est orienté vers le bas, la longueur du ressort s'identifie donc directement à la coordonnée $y$. 

La force de rappel s'écrit $\vv{F}=-k(y-\ell_0)\vv{e_y}$. On en déduit donc (en calculant le travail élémentaire ou par intégration directe) que 
\[
E_{pe}(y)=\frac{1}{2}k(y-\ell_0)^2 + \mathrm{C^{te}}.
\]

Or, on veut $E_{pe}(y=0)=0$, d'où $\mathrm{C^{te}}=-\frac{1}{2}k{\ell_0}^2$. Ainsi, on a
\[
E_{pe}(y) = \frac{1}{2}k(y-\ell_0)^2 -\frac{1}{2}k{\ell_0}^2.
\]
\end{corrige}

% ************************************** %


% ^^^^^^^^^^^^^^^^^^^^^^^^^^^^^^^^^^^^^^ %
%               Question                 %
% ^^^^^^^^^^^^^^^^^^^^^^^^^^^^^^^^^^^^^^ %

\begin{enonce}
$E_{pe}(x) = $
\end{enonce}

\reponse{\small$\frac{1}{2}k\left(\frac{x}{\cos(\beta)}-\ell_0\right)^2-\frac{1}{2}k\left(\frac{L}{\sin(\beta)}-\ell_0\right)^2$}

\begin{corrigeNewpage}
On calcule d'abord la longueur $\ell$ du ressort en fonction de la coordonnée $x$. Un peu de trigonométrie donne $\cos(\beta) = \frac{x}{\ell}$ d'où $\ell = \frac{x}{\cos(\beta)}$. Par rapport à la coordonnée $\ell$ (mesurée le long de l'axe $(\ptO\ptA)$), l'énergie potentielle vaut donc :
\[
E_{pe}(\ell) = \frac{1}{2}k(\ell-\ell_0)^2+\mathrm{C^{te}}.
\]
On a donc
\[
E_{pe}(x) = \frac{1}{2}k\left(\frac{x}{\cos(\beta)}-\ell_0\right)^2+\mathrm{C^{te}}.
\]

On détermine alors la constante afin d'avoir $E_{pe}(\ptA)=0$. Lorsque le point M est en A, la longueur du ressort vaut $\ell(\ptA)=\frac{L}{\sin(\beta)}$. On résout donc : \minusMedskip
\[
E_{pe}(\ell(\ptA)) = \frac{1}{2}k\left(\frac{L}{\sin(\beta)}-\ell_0\right)^2+\mathrm{C^{te}} = 0 
\quad\text{ce qui donne}\quad
\mathrm{C^{te}}=-\frac{1}{2}k\left(\frac{L}{\sin(\beta)}-\ell_0\right)^2.
\]
Finalement, on trouve
\[
E_{pe}(x)=\frac{1}{2}k\left(\frac{x}{\cos(\beta)}-\ell_0\right)^2-\frac{1}{2}k\left(\frac{L}{\sin(\beta)}-\ell_0\right)^2.
\]
\end{corrigeNewpage}

% ************************************** %


% ^^^^^^^^^^^^^^^^^^^^^^^^^^^^^^^^^^^^^^ %
%               Question                 %
% ^^^^^^^^^^^^^^^^^^^^^^^^^^^^^^^^^^^^^^ %

\begin{enonce}
$E_{pe}(x) = $
\end{enonce}

\reponse{$E_0+k(x-\ell_0)^2$}

\begin{corrige}
La masse centrale est soumise aux forces de rappel des deux ressorts :
\begin{itemize}
\item La longueur du ressort de gauche vaut $x$. La force exercée par celui-ci sur la masse s'exprime donc comme $\vv{F_g}=-k(x-\ell_0)\vv{e_x}$, d'où une énergie potentielle (à une constante près) $E_{p,g} = \frac{1}{2}k(x-\ell_0)^2$.

\item La longueur du ressort de droite vaut $2\ell_0-x$. La force exercée par celui-ci sur la masse s'exprime donc comme $\vv{F_d}=k(2\ell_0-x-\ell_0)\vv{e_x} = k(\ell_0-x)\vv{e_x}$ (attention au signe devant $k$ qui doit être cohérent), d'où une énergie potentielle (à une constante près) $E_{p,d} = \frac{1}{2}k(\ell_0-x)^2$.
\end{itemize}
En additionnant les deux contributions, et en demandant que $E_{pe}(\ell_0)=E_0$, on obtient alors $E_{pe}(x) = E_0+k(x-\ell_0)^2$.
\end{corrige}

% ************************************** %

% =============================================================================== %
\finEntrainement
% =============================================================================== %



\newpage

% •••••••••••••••••••••••••••••••••••••••••••••••••••••••••••••••••••••••••••• %
% •••••••••••••••••••••••••••••••••••••••••••••••••••••••••••••••••••••••••••• %
\sectionFicheEntrainement{Travail d'une force}
% •••••••••••••••••••••••••••••••••••••••••••••••••••••••••••••••••••••••••••• %
% •••••••••••••••••••••••••••••••••••••••••••••••••••••••••••••••••••••••••••• %







% ********************************************************************* %
%                           ENTRAÎNEMENT                                % 
% ********************************************************************* %
% ================ Métadonnées sur l'entraînement ===================== %

\titreEntrainementFacultatif	{Une force de frottement}
\hauteurLargeurCadreReponse		{7mm}{3cm}
\dureeResolutionFacultative		{3} % 1, 2, 3 ou 4
\typeEntrainement				{Newton} % Newton ou QCM ou AN ou vide (exercice "normal")
\nombreColonnesQuestions		{1} % vide, 1, 2, 3, etc.
\avecPlusieursQuestions			{Y} % Y ou N
\initialisationEntrainement
% ===================================================================== %

On considère le travail $W_{\ptA \ptB}=\int_{\ptA}^{\ptB} \vv{F}\cdot \vv{ \d{}\ell}$ d'une force de frottement $\vv{F} = -h\frac{\vv{v}}{\|\vv{v}\|}$ où $\vv{v}$ est le vecteur vitesse du point matériel subissant la force et $h$ est une constante. 

\medskip

Déterminer $W$ pour les chemins suivants : 
% =============================================================================== %
\debutEntrainement
% =============================================================================== %

% ^^^^^^^^^^^^^^^^^^^^^^^^^^^^^^^^^^^^^^ %
%               Question                 %
% ^^^^^^^^^^^^^^^^^^^^^^^^^^^^^^^^^^^^^^ %

\begin{enonce}
Un segment reliant $\ptA(0,0)$ et $\ptB(\ell,0)$
\end{enonce}

\reponse{$-h\ell$}

\begin{corrige}
Déterminons le travail élémentaire : on a
\[
\delta W = \vv{F}\cdot\d{}\vv{\ptO{}\ptM} = -\frac{h}{||\vv{v}||}\vv{v}\cdot\d{}\vv{\ptO{}\ptM}.
\]
Or, par construction, les vecteurs vitesse et déplacement élémentaire sont colinéaires, d'où :
\[
\delta W = -h \d \ptO\ptM.
\]
Par intégration, on a donc 
\[
W = \int_{\ptA\ptB}{-h \d \ptO\ptM} = -h \int_{\ptA\ptB}{\d \ptO\ptM}=-h\ell.
\]
Les autres cas se calculent semblablement.
\end{corrige}

% ************************************** %


% ^^^^^^^^^^^^^^^^^^^^^^^^^^^^^^^^^^^^^^ %
%               Question                 %
% ^^^^^^^^^^^^^^^^^^^^^^^^^^^^^^^^^^^^^^ %

\begin{enonce}
Un arc de cercle d'angle $\alpha$ et de rayon $R$
\end{enonce}

\reponse{$-hR\alpha$}

% ************************************** %


% ^^^^^^^^^^^^^^^^^^^^^^^^^^^^^^^^^^^^^^ %
%               Question                 %
% ^^^^^^^^^^^^^^^^^^^^^^^^^^^^^^^^^^^^^^ %

\begin{enonce}
Un rectangle $\ptA\ptB\ptC\ptD$ de côtés $a$ et $b$
\end{enonce}

\reponse{$-(2a+2b)h$}

% ************************************** %


% ^^^^^^^^^^^^^^^^^^^^^^^^^^^^^^^^^^^^^^ %
%               Question                 %
% ^^^^^^^^^^^^^^^^^^^^^^^^^^^^^^^^^^^^^^ %

\begin{enonce}
Un triangle $\ptA\ptB\ptC$ de côtés $a,b,c$
\end{enonce}

\reponse{$-(a+b+c)h$}

% ************************************** %


% ^^^^^^^^^^^^^^^^^^^^^^^^^^^^^^^^^^^^^^ %
%               Question                 %
% ^^^^^^^^^^^^^^^^^^^^^^^^^^^^^^^^^^^^^^ %

\begin{enonce}
En comparant les résultats obtenus, peut-on dire que la force est conservative ?

	\begin{listeQCM2Colonnes}
	\item Oui
	\item Non
	\end{listeQCM2Colonnes}
\end{enonce}

\reponse{\reponseB{}}

\begin{corrige}
Si la force était conservative, son travail ne dépendrait que des points de départ et d'arrivée, et serait donc nul sur un chemin fermé (points de départ et d'arrivée confondus). Ce n'est pas le cas pour les chemin c) et d), la force n'est donc pas conservative.
\end{corrige}

% ************************************** %


% =============================================================================== %
\finEntrainement
% =============================================================================== %







% •••••••••••••••••••••••••••••••••••••••••••••••••••••••••••••••••••••••••••• %
% •••••••••••••••••••••••••••••••••••••••••••••••••••••••••••••••••••••••••••• %
\sectionFicheEntrainement{Théorèmes énergétiques}
% •••••••••••••••••••••••••••••••••••••••••••••••••••••••••••••••••••••••••••• %
% •••••••••••••••••••••••••••••••••••••••••••••••••••••••••••••••••••••••••••• %


% ********************************************************************* %
%                           ENTRAÎNEMENT                                % 
% ********************************************************************* %
% ================ Métadonnées sur l'entraînement ===================== %

\titreEntrainementFacultatif	{Freinage et variation d'énergie cinétique}
\hauteurLargeurCadreReponse		{8mm}{1cm}
\dureeResolutionFacultative		{2} % 1, 2, 3 ou 4
\typeEntrainement				{Newton} % Newton ou QCM ou AN ou vide (exercice "normal")
\nombreColonnesQuestions		{1} % vide, 1, 2, 3, etc.
\avecPlusieursQuestions			{N} % Y ou N
\initialisationEntrainement
% ===================================================================== %

On considère une voiture (assimilée à un point matériel de masse $m$) se déplaçant le long d'une route rectiligne horizontale et 
dont la vitesse initiale au début de la phase de freinage vaut
$\vv{v}=v_0\vv{e_x}$. 

En freinant, le véhicule est soumis à une force de frottement $\vv{F}=-h \vv{e_x}$.

% =============================================================================== %
\debutEntrainement
% =============================================================================== %

% ^^^^^^^^^^^^^^^^^^^^^^^^^^^^^^^^^^^^^^ %
%               Question                 %
% ^^^^^^^^^^^^^^^^^^^^^^^^^^^^^^^^^^^^^^ %

\begin{enonce}
	Quelle est l'expression de la distance d'arrêt $d$ de la voiture ?
	
	\begin{listeQCM3Colonnes}
	\item $\frac{2m{v_0}^2}{h}$
	\item $\frac{m{v_0}^2}{h}$
	\item $\frac{m{v_0}^2}{2h}$
	\end{listeQCM3Colonnes}
		
\end{enonce}

\reponse{\reponseC{}}

\begin{corrige}
On applique le théorème de l'énergie cinétique entre le point de départ et le point d'arrêt. L'entraînement précédent permet d'affirmer que le travail de la force de frottement vaut $-hd$. On a donc :
\[
\Delta E_c = 0-\frac{1}{2}m{v_0}^2 = -hd \qquad\text{donc}\qquad d = \frac{m{v_0}^2}{2h}.
\]
\end{corrige}

% ************************************** %

% =============================================================================== %
\finEntrainement
% =============================================================================== %






% ********************************************************************* %
%                           ENTRAÎNEMENT                                % 
% ********************************************************************* %
% ================ Métadonnées sur l'entraînement ===================== %

\titreEntrainementFacultatif	{Pendule simple}
\hauteurLargeurCadreReponse		{8mm}{5cm}
\dureeResolutionFacultative		{2} % 1, 2, 3 ou 4
\typeEntrainement				{} % Newton ou QCM ou AN ou vide (exercice "normal")
\nombreColonnesQuestions		{1} % vide, 1, 2, 3, etc.
\avecPlusieursQuestions			{Y} % Y ou N
\initialisationEntrainement
% ===================================================================== %

Un pendule simple est constitué d'un fil de longueur $\ell=\np[m]{1,0}$ auquel est accroché une masse $m=\np[g]{100}$. 

À $t=0$, on donne à cette masse une vitesse horizontale $\vv{v_0}=v_0  \vv{e_x}$ où $v_0=\np[m\cdot s^{-1}]{2,0}$. 

On note $\theta_0$ l'angle pour lequel la masse rebrousse chemin.


% =============================================================================== %
\debutEntrainement
% =============================================================================== %

% ^^^^^^^^^^^^^^^^^^^^^^^^^^^^^^^^^^^^^^ %
%               Question                 %
% ^^^^^^^^^^^^^^^^^^^^^^^^^^^^^^^^^^^^^^ %

\begin{enonce}
Exprimer $\cos(\theta_0)$
\end{enonce}

\reponse{$1-\frac{{v_0}^2}{2g\ell}$}

\begin{corrige}
La masse n'est soumise qu'au poids, force conservative, et à la tension du fil qui ne travaille pas car elle reste orthogonale au mouvement. L'énergie mécanique se conserve donc entre le point de départ et le point de rebroussement.
\begin{itemize}
\item Au départ, $E_m = E_c = \frac{1}{2}m{v_0}^2$ (on pose $z=0$ pour la position initiale de la masse, et on prend $E_p(0)=0$)

\item Au moment du rebroussement, $E_m = E_p = mgz(\theta_0) =  mg\ell(1-\cos(\theta_0))$, car on a alors $z(\theta)=\ell-\ell\cos(\theta)$.
\end{itemize}

Ainsi, on a 
\[
\frac{1}{2}m{v_0}^2 = mg\ell(1-\cos(\theta_0)) \qquad\text{donc}\qquad \cos(\theta_0) = 1-\frac{{v_0}^2}{2g\ell}.
\]
\end{corrige}

% ************************************** %


% ^^^^^^^^^^^^^^^^^^^^^^^^^^^^^^^^^^^^^^ %
%               Question                 %
% ^^^^^^^^^^^^^^^^^^^^^^^^^^^^^^^^^^^^^^ %

\begin{enonce}
Calculer $\theta_0$
\end{enonce}

\reponse{$\np[rad]{0,65}=\ang{37}$}

% ************************************** %

% =============================================================================== %
\finEntrainement
% =============================================================================== %


\newpage






% ********************************************************************* %
%                           ENTRAÎNEMENT                                % 
% ********************************************************************* %
% ================ Métadonnées sur l'entraînement ===================== %

\titreEntrainementFacultatif	{Trampoline simplifié}
\hauteurLargeurCadreReponse		{8mm}{5cm}
\dureeResolutionFacultative		{3} % 1, 2, 3 ou 4
\typeEntrainement				{AN} % Newton ou QCM ou AN ou vide (exercice "normal")
\nombreColonnesQuestions		{1} % vide, 1, 2, 3, etc.
\avecPlusieursQuestions			{Y} % Y ou N
\initialisationEntrainement
% ===================================================================== %

Un ressort de longueur à vide $\ell_0=\np[cm]{30}$, de raideur $k=\np[N\cdot m^{-1}]{1,0e3}$, sans masse, est posé sur le sol à la verticale. On lâche d'une hauteur $H=\np[m]{2,0}$ et sans vitesse initiale une masse ponctuelle $m=\np[kg]{1,0}$. Après une durée de chute libre sans frottement, la masse atteint le ressort, le comprime jusqu'à ce que celui-ci la propulse vers le haut comme le ferait un trampoline. 

En admettant que la masse quitte le ressort quand $z=\ell_0$, calculer :

% =============================================================================== %
\debutEntrainement
% =============================================================================== %

% ^^^^^^^^^^^^^^^^^^^^^^^^^^^^^^^^^^^^^^ %
%               Question                 %
% ^^^^^^^^^^^^^^^^^^^^^^^^^^^^^^^^^^^^^^ %

\begin{enonce}
La vitesse de la masse lors du contact avec le ressort 
\end{enonce}

\reponse{$\SI{5,8}{m\cdot s^{-1}}$}

\begin{corrige}
En appliquant le théorème de l'énergie mécanique entre le début et la fin de la chute libre on a :
\[
E_m(t_\mathrm{fin \; chute})- E_m(t_\mathrm{début \; chute}) =  \frac{1}{2} m {v_0}^2 - mg (H-\ell_0).
\]
Les forces étant conservatives, l'énergie mécanique est conservée et on a donc 
\[
v_0 = \sqrt{2g (H-\ell_0)}=\sqrt{2 \times \np[m\cdot s^{-2}]{9,81} \times (\np[m]{2,0}-\np[m]{0,30})}=\SI{5,8}{m\cdot s^{-1}}.
\]
\end{corrige}

% ************************************** %


% ^^^^^^^^^^^^^^^^^^^^^^^^^^^^^^^^^^^^^^ %
%               Question                 %
% ^^^^^^^^^^^^^^^^^^^^^^^^^^^^^^^^^^^^^^ %

\begin{enonce}
L'altitude minimale atteinte par la masse
\end{enonce}

\reponse{\np[m]{0,11}}

\begin{corrige}
La masse n'est soumise qu'à des forces conservatives : son poids, ainsi que la force de rappel du ressort. On peut donc appliquer la conservation de l'énergie mécanique entre la position d'arrivée sur le ressort $z=\ell_0$, et la position d'altitude minimale $z=z_m$ (pour laquelle la vitesse s'annule). On a donc
\[
\frac{1}{2}m{v_0}^2+mg\ell_0 = mgz_m+\frac{1}{2}k(z_m-\ell_0)^2.
\]
Ainsi, après calcul, on trouve $\frac{1}{2}kz_m^2+(mg-k\ell_0)z_m+\frac{1}{2}k{\ell_0}^2-\frac{1}{2}m{v_0}^2-mg\ell_0=0$.

On ne demande qu'une réponse numérique, on peut donc passer aux valeurs numériques pour simplifier la résolution :
\[
500 z_m^2-\np{290,2} z_m+\np{25,4} = 0.
\]
Cette équation possède deux solutions, $z_1\approx\np[m]{0,47}$ et $z_2\approx\np[m]{0,11}$. La première solution correspond à une position supérieure en altitude à la position initiale, et n'est donc pas celle qui nous intéresse. On retient donc
$z_m = \np[m]{0,11}$.
\end{corrige}

% ************************************** %


% ^^^^^^^^^^^^^^^^^^^^^^^^^^^^^^^^^^^^^^ %
%               Question                 %
% ^^^^^^^^^^^^^^^^^^^^^^^^^^^^^^^^^^^^^^ %

\begin{enonce}
L'altitude maximale de la masse (en fin de remontée)
\end{enonce}

\reponse{$\np[m]{2,0}$}

\begin{corrige}
La masse n'étant soumise qu'à des forces conservatives, elle revient en $x=\ell_0$ avec la même vitesse qu'elle avait en arrivant, à savoir $v_0$. Elle atteint donc une altitude maximale quand sa vitesse s'annule en $z=H$.
\end{corrige}

% ************************************** %

% =============================================================================== %
\finEntrainement
% =============================================================================== %







% ********************************************************************* %
%                           ENTRAÎNEMENT                                % 
% ********************************************************************* %
% ================ Métadonnées sur l'entraînement ===================== %

\titreEntrainementFacultatif	{Oscillateur vertical}
\hauteurLargeurCadreReponse		{10mm}{9cm}
\dureeResolutionFacultative		{3}% 1, 2, 3 ou 4
\typeEntrainement				{Newton} % Newton ou QCM ou AN ou vide (exercice "normal")
\nombreColonnesQuestions		{1} % vide, 1, 2, 3, etc.
\avecPlusieursQuestions			{Y} % Y ou N
\initialisationEntrainement
% ===================================================================== %

Un point \ptM  de masse $m$ est accroché à une paroi horizontale fixe par l'intermédiaire d'un ressort de raideur $k$ et de longueur à vide $\ell_0$. Son mouvement s'effectue dans un liquide qui produit une force de frottements fluides linéaire $\vv{F}
=-\alpha\vv{v}$, où $\alpha>0$. On néglige la poussée d'Archimède, on ne considère que des mouvements verticaux dans le champ de pesanteur $\vv{g}$.


\begin{center}
\subimport{_images/}{Oscillateur_GBU_v2.tex}
\end{center}


% =============================================================================== %
\debutEntrainement
% =============================================================================== %

% ^^^^^^^^^^^^^^^^^^^^^^^^^^^^^^^^^^^^^^ %
%               Question                 %
% ^^^^^^^^^^^^^^^^^^^^^^^^^^^^^^^^^^^^^^ %

\begin{enonce}
On note $z$ la position de $\ptM$ par rapport à $\ptO$. 

Déterminer, par une méthode énergétique, l'équation différentielle vérifiée par $z$.

\end{enonce}

\reponse{$\ddot z+\frac{\alpha}{m}\dot z + \frac{k}{m} z = g+\frac{k\ell_0}{m}$}

\begin{corrige}
On choisit un axe vertical descendant de manière à pouvoir identifier $z$ à la distance $\ptO\ptM$, qui est la longueur du ressort.

Afin de déterminer l'équation différentielle, on souhaite appliquer le théorème de la puissance cinétique. Or, 
\begin{itemize}
\item[$\bullet$] la puissance du poids vaut $m\vv{g}\cdot\vv{v}=mg\dot z$ (axe descendant) ;
\item[$\bullet$] la puissance de la force de rappel vaut $-k(z-\ell_0)\vv{e_z}\cdot\vv{v} = -k(z-\ell_0)\dot z$ ;
\item[$\bullet$] la puissance de la force de frottements fluides vaut $-\alpha\vv{v}\cdot\vv{v}=-\alpha\dot z^2$.
\end{itemize}

Le théorème de la puissance cinétique donne alors :
\[
\diff{E_c}{t} = \diff{}{t}\left(\frac{1}{2}m\dot z^2\right) = m\dot z \ddot z = 
mg\dot z -k(z-\ell_0)\dot z -\alpha\dot z^2.
\]
D'où finalement : $\ddot z +\frac{\alpha}{m}\dot z + \frac{k}{m}z =  g+\frac{k \ell_0}{m}$.
\end{corrige}

% ************************************** %


% ^^^^^^^^^^^^^^^^^^^^^^^^^^^^^^^^^^^^^^ %
%               Question                 %
% ^^^^^^^^^^^^^^^^^^^^^^^^^^^^^^^^^^^^^^ %

\begin{enonce}
On note à présent $\zeta$ la position de $\ptM$ par rapport à sa position à l'équilibre. 

Déterminer l'équation différentielle vérifiée par $\zeta$.

\end{enonce}

\reponse{$\zeta +\frac{\alpha}{m}\dot \zeta + \frac{k}{m}\zeta = 0$}


\begin{corrigeNewpage}
On détermine la position d'équilibre en projetant la première loi de Newton sur l'axe vertical descendant : \vspace{-1mm}
\[
mg-k(z_{\textrm{eq}}-\ell_0) = 0 \quad\text{donc}\quad z_{\textrm{eq}} = \ell_0+\frac{mg}{k}.
\]
On obtient $z_{\textrm{eq}}>\ell_0$, ce qui est physiquement cohérent.

On pose donc $\zeta = z-z_{\textrm{eq}}$. En réinjectant dans l'équation différentielle obtenue précédemment, on obtient :
\[
\ddot \zeta +\frac{\alpha}{m}\dot \zeta + \frac{k}{m}\left(\zeta+\ell_0+\frac{mg}{k}\right) = g+\frac{k \ell_0}{m} 
\qquad\text{donc}\qquad
\ddot \zeta +\frac{\alpha}{m}\dot \zeta + \frac{k}{m}\zeta = 0.
\]

{\itshape On peut également obtenir cette équation en écrivant la force de rappel par rapport à la variable $\zeta$ et en en déduisant l'énergie potentielle associée.}
\end{corrigeNewpage}

% ************************************** %

% =============================================================================== %
\finEntrainement
% =============================================================================== %




\newpage


% •••••••••••••••••••••••••••••••••••••••••••••••••••••••••••••••••••••••••••• %
% •••••••••••••••••••••••••••••••••••••••••••••••••••••••••••••••••••••••••••• %
\sectionFicheEntrainement{Mouvements conservatifs et positions d'équilibre}
% •••••••••••••••••••••••••••••••••••••••••••••••••••••••••••••••••••••••••••• %
% •••••••••••••••••••••••••••••••••••••••••••••••••••••••••••••••••••••••••••• %



% ********************************************************************* %
%                           ENTRAÎNEMENT                                % 
% ********************************************************************* %
% ================ Métadonnées sur l'entraînement ===================== %

\titreEntrainementFacultatif	{Profils d'énergies potentielles}
\hauteurLargeurCadreReponse		{7mm}{1.5cm}
\dureeResolutionFacultative		{2} % 1, 2, 3 ou 4
\typeEntrainement				{QCM} % Newton ou QCM ou AN ou vide (exercice "normal")
\nombreColonnesQuestions		{2} % vide, 1, 2, 3, etc.
\avecPlusieursQuestions			{Y} % Y ou N
\initialisationEntrainement
% ===================================================================== %

Les quatre profils suivants représentent la fonction énergie potentielle 
\[
E_p(x) = \frac{\alpha}{x}+\frac{\beta}{x^2}\minusSmallskip
\]
avec $\alpha,\beta$ des réels non nuls.


\begin{center}
\subimport{_images/}{Profil_potentiel_4_CBD_v2.tex}
\end{center}


Attribuer à chacune des figures ci-dessus les bons signes pour $\alpha$ et $\beta$, en indiquant laquelle des réponses suivantes est la bonne :
	\begin{listeQCM2Colonnes}
	\item $\alpha>0$ \ {}et \  $\beta>0$
	\item $\alpha>0$ \ {}et \  $ \beta<0$
	\item $\alpha<0$ \ {}et \  $ \beta>0$
	\item $\alpha<0$ \ {}et \  $ \beta<0$
	\end{listeQCM2Colonnes}

\bigskip

% =============================================================================== %
\debutEntrainement
% =============================================================================== %

% ^^^^^^^^^^^^^^^^^^^^^^^^^^^^^^^^^^^^^^ %
%               Question                 %
% ^^^^^^^^^^^^^^^^^^^^^^^^^^^^^^^^^^^^^^ %

\begin{enonce}
	Énergie potentielle \no{}1
\end{enonce}

\reponse{\reponseB{}}

\begin{corrige}
Au voisinage de $x=0^+$, la fonction énergie potentielle est équivalente à $\beta/x^2$. Ici, la fonction représentée par le graphe tend vers $-\infty$ en $0$, on a nécessairement $\beta<0$.

Pour $x\to+\infty$, la fonction énergie potentielle est équivalente à $\alpha/x$. Ici, la fonction représentée par le graphe tend vers $0^+$ en $+\infty$, on a nécessairement $\alpha>0$.

{\itshape Ce potentiel est physiquement impossible car $E_p(x\to 0^+) \longto -\infty$ : l'énergie potentielle n'est pas bornée inférieurement, on pourrait donc théoriquement utiliser ce potentiel pour extraire une quantité infinie d'énergie.}
\end{corrige}

% ************************************** %


% ^^^^^^^^^^^^^^^^^^^^^^^^^^^^^^^^^^^^^^ %
%               Question                 %
% ^^^^^^^^^^^^^^^^^^^^^^^^^^^^^^^^^^^^^^ %

\begin{enonce}
	Énergie potentielle \no{}2
\end{enonce}

\reponse{\reponseD{}}

% ************************************** %


% ^^^^^^^^^^^^^^^^^^^^^^^^^^^^^^^^^^^^^^ %
%               Question                 %
% ^^^^^^^^^^^^^^^^^^^^^^^^^^^^^^^^^^^^^^ %

\begin{enonce}
	Énergie potentielle \no{}3	
\end{enonce}

\reponse{\reponseA{}}

% ************************************** %


% ^^^^^^^^^^^^^^^^^^^^^^^^^^^^^^^^^^^^^^ %
%               Question                 %
% ^^^^^^^^^^^^^^^^^^^^^^^^^^^^^^^^^^^^^^ %

\begin{enonce}
	Énergie potentielle \no{}4	
\end{enonce}

\reponse{\reponseC{}}

% ************************************** %


% =============================================================================== %
\finEntrainement
% =============================================================================== %






% ********************************************************************* %
%                           ENTRAÎNEMENT                                % 
% ********************************************************************* %
% ================ Métadonnées sur l'entraînement ===================== %

\titreEntrainementFacultatif	{Autour d'une position d'équilibre}
\hauteurLargeurCadreReponse		{8mm}{4cm}
\dureeResolutionFacultative		{3} % 1, 2, 3 ou 4
\typeEntrainement				{} % Newton ou QCM ou AN ou vide (exercice "normal")
\basiqueEtTransversal			{Y}
\nombreColonnesQuestions		{1} % vide, 1, 2, 3, etc.
\avecPlusieursQuestions			{Y} % Y ou N
\initialisationEntrainement
% ===================================================================== %

On donne l'expression de potentiels $E_p$ dans chacun desquels évolue un point matériel de masse $m$. 
\medskip

Déterminer dans chaque cas la position d'équilibre stable.
\smallskip

% =============================================================================== %
\debutEntrainement
% =============================================================================== %

% ^^^^^^^^^^^^^^^^^^^^^^^^^^^^^^^^^^^^^^ %
%               Question                 %
% ^^^^^^^^^^^^^^^^^^^^^^^^^^^^^^^^^^^^^^ %



\begin{enonce}
Pour $E_p(\theta)=mg\ell(1-\cos(\theta))$ : \minusMedskip \\
$\theta_{\textrm{eq}} = $
\end{enonce}

\reponse{$0$}

\begin{corrige}
La position d'équilibre stable correspond à l'état qui minimise l'énergie potentielle.

\begin{itemize}
\item
Déterminons le minimum de l'énergie potentielle $E_p(\theta)=mg\ell(1-\cos(\theta))$ en cherchant la valeur $\theta_\text{eq}$ telle que 
\[
\diff{E_p}{\theta}(\theta_\text{eq})=0
\qquad\text{et}\qquad
\diff[2]{E_p}{\theta}(\theta_\text{eq})>0.
\]
La première égalité donne $
\diff{E_p}{\theta}(\theta_\text{eq})=
mg\ell\sin(\theta_\text{eq})=0$ et donc $\theta_\text{eq} \equiv 0\ [\pi]$.

Finalement, en tenant compte de $\diff[2]{E_p}{\theta}(\theta_\text{eq})>0$, on trouve $\theta_\text{eq} \equiv 0\ [2\pi]$.

\item On aurait pu remarquer que les minima de $mg\ell(1-\cos(\theta))$ correspondent aux maxima de $\cos(\theta)$, qui sont bien les $\theta_\text{eq}\equiv0\ [2\pi]$.
\end{itemize}
\end{corrige}

% ************************************** %




% ^^^^^^^^^^^^^^^^^^^^^^^^^^^^^^^^^^^^^^ %
%               Question                 %
% ^^^^^^^^^^^^^^^^^^^^^^^^^^^^^^^^^^^^^^ %

\begin{enonce}
Pour $E_p(z)=\frac{1}{2}\kappa z^2+\frac{1}{4}\lambda z^4$ \quad avec $\kappa>0$ et $\lambda<0$ : \minusBigskip\\
$z_{\textrm{eq}} = $
\end{enonce}

\reponse{$0$}

\begin{corrige}
On dérive l'énergie potentielle, en écrivant : \minusSmallskip
\[
\diff{E_p}{z} = \kappa z + \lambda z^3.\minusSmallskip
\]
L'équation $\diff{E_p}{z} = 0$ a alors trois solutions : 
$z_1=0$,  
$z_2=\sqrt{-\frac{\kappa}{\lambda}}$
et 
$z_3=-\sqrt{-\frac{\kappa}{\lambda}}$.

Il s'agit des positions d'équilibre de ce potentiel.

On dérive une seconde fois afin d'étudier la stabilité. On a 
$\diff[2]{E_p}{z} = \kappa+3\lambda z^2$.

Finalement, on obtient : \minusMedskip
\begin{align*}
\diff[2]{E_p}{z}(z=z_1)&=\kappa >0 \\
 \diff[2]{E_p}{z}(z=z_2)&=\kappa+3\lambda\left(-\frac{\kappa}{\lambda}\right)=-2\kappa <0 \\
 \diff[2]{E_p}{z}(z=z_3)&=\kappa+3\lambda\left(-\frac{\kappa}{\lambda}\right)=-2\kappa <0.
\end{align*}
Seule $z_1=0$ est une position d'équilibre stable.
\end{corrige}

% ************************************** %




% ^^^^^^^^^^^^^^^^^^^^^^^^^^^^^^^^^^^^^^ %
%               Question                 %
% ^^^^^^^^^^^^^^^^^^^^^^^^^^^^^^^^^^^^^^ %

\begin{enonce}
Pour $E_p(x) = U_0\,\eEuler^{\beta x^2}$ avec  $U_0,\beta >0$ : \minusMedskip \\
$x_{\textrm{eq}} = $
\end{enonce}

\reponse{$0$}

\begin{corrige}
On calcule la dérivée de l'énergie potentielle : \minusSmallskip
\[
\diff{E_p}{x} = 2U_0\beta x \eEuler^{\beta x^2} \minusSmallskip
\]
qui montre que $\diff{E_p}{x}$ s'annule pour $x=0$, qui est donc une position d'équilibre. 

Pour étudier sa stabilité, on dérive une seconde fois :
\[
\diff{^2E_p}{x^2}=2U_0\beta(1+2\beta x^2)\eEuler^{\beta x^2}
\]
qui en $x=0$ vaut $2U_0\beta>0$. Cette position d'équilibre est donc bien stable.
\end{corrige}

% ************************************** %




% ^^^^^^^^^^^^^^^^^^^^^^^^^^^^^^^^^^^^^^ %
%               Question                 %
% ^^^^^^^^^^^^^^^^^^^^^^^^^^^^^^^^^^^^^^ %

\begin{enonce}
Pour $E_p(\phi) = E_0\sin^2(\phi-a)$\quad  avec $E_0>0$, $\phi\in[0,\pi[$ et  $a\in\left[0,\frac{\pi}{2}\right]$ : \minusBigskip \\
$\phi_{\textrm{eq}} = $
\end{enonce}

\reponse{$a$}

\begin{corrige}
On calcule la dérivée de l'énergie potentielle :
\[
\diff{E_p}{\phi} =2E_0\sin(\phi-a)\cos(\phi-a).
\]
Ainsi, $\diff{E_p}{\phi}$ s'annule pour $\phi=a$ et $\phi=a+\frac{\pi}{2}$, qui sont les positions d'équilibre dans l'intervalle $[0,\pi[$. 

Afin d'étudier leur stabilité, on dérive une seconde fois :
\[
\diff{^2E_p}{\phi^2}=2E_0(\cos^2(\phi-a)-\sin^2(\phi-a)).
\]

\begin{itemize}
\item
On calcule ensuite $\diff{^2E_p}{\phi^2}(\phi=a) = 2E_0$.
Ce dernier terme étant positif, la position d'équilibre $\phi=a$ est donc stable.

\item
Pour l'autre position d'équilibre, on a $\diff{^2E_p}{\phi^2}(\phi=a+\pi/2) = -2E_0$.
Cette dérivée seconde étant négative, la position d'équilibre $\phi=a+\pi/2$ est instable.
\end{itemize}
\end{corrige}

% ************************************** %

% =============================================================================== %
\finEntrainement
% =============================================================================== %




% ********************************************************************* %
%                           ENTRAÎNEMENT                                % 
% ********************************************************************* %
% ================ Métadonnées sur l'entraînement ===================== %

\pointInterrogation				{Y}
\titreEntrainementFacultatif	{État lié ou état de diffusion}
\hauteurLargeurCadreReponse		{6mm}{2cm}
\dureeResolutionFacultative		{1} % 1, 2, 3 ou 4
\typeEntrainement				{Newton} % Newton ou QCM ou AN ou vide (exercice "normal")
\nombreColonnesQuestions		{2} % vide, 1, 2, 3, etc.
\avecPlusieursQuestions			{Y} % Y ou N
\initialisationEntrainement
% ===================================================================== %

On considère le profil suivant d'énergie potentielle (les abscisses étoilées et l'abscisse $x_3$ serviront dans l'entraînement suivant). 

Pour chaque état suivant, étant donné les valeurs de l'énergie mécanique et de la position initiale d'un point matériel, dire si ce dernier se trouve :
	\begin{listeQCM2Colonnes}
	\item dans un état lié
	\item un état de diffusion.
	\end{listeQCM2Colonnes}

\smallskip

\begin{center}
\subimport{_images/}{Profil_potentiel_CBD_v3.tex}
\end{center}

% =============================================================================== %
\debutEntrainement
% =============================================================================== %

% ^^^^^^^^^^^^^^^^^^^^^^^^^^^^^^^^^^^^^^ %
%               Question                 %
% ^^^^^^^^^^^^^^^^^^^^^^^^^^^^^^^^^^^^^^ %

\begin{enonce}
$E_m=E_1$ \ et \ $ x(0)=x_1$
\end{enonce}

\reponse{\reponseA{}}

% ************************************** %


% ^^^^^^^^^^^^^^^^^^^^^^^^^^^^^^^^^^^^^^ %
%               Question                 %
% ^^^^^^^^^^^^^^^^^^^^^^^^^^^^^^^^^^^^^^ %

\begin{enonce}
$E_m=E_1$ \ et \ $ x(0)=x_2$
\end{enonce}

\reponse{\reponseA{}}

% ************************************** %


% ^^^^^^^^^^^^^^^^^^^^^^^^^^^^^^^^^^^^^^ %
%               Question                 %
% ^^^^^^^^^^^^^^^^^^^^^^^^^^^^^^^^^^^^^^ %

\begin{enonce}
$E_m=E_2$ \ et \ $ x(0)=x_1$
\end{enonce}

\reponse{\reponseA{}}

% ************************************** %


% ^^^^^^^^^^^^^^^^^^^^^^^^^^^^^^^^^^^^^^ %
%               Question                 %
% ^^^^^^^^^^^^^^^^^^^^^^^^^^^^^^^^^^^^^^ %

\begin{enonce}
$E_m=E_2$ \ et \ $ x(0)=x_2$
\end{enonce}

\reponse{\reponseB{}}

% ************************************** %


% ^^^^^^^^^^^^^^^^^^^^^^^^^^^^^^^^^^^^^^ %
%               Question                 %
% ^^^^^^^^^^^^^^^^^^^^^^^^^^^^^^^^^^^^^^ %

\begin{enonce}
$E_m=E_3$ \ et \ $ x(0)=x_1$
\end{enonce}

\reponse{\reponseB{}}

% ************************************** %


% ^^^^^^^^^^^^^^^^^^^^^^^^^^^^^^^^^^^^^^ %
%               Question                 %
% ^^^^^^^^^^^^^^^^^^^^^^^^^^^^^^^^^^^^^^ %

\begin{enonce}
$E_m=E_3$ \ et \ $ x(0)=x_2$
\end{enonce}

\reponse{\reponseB{}}

% ************************************** %

% =============================================================================== %
\finEntrainement
% =============================================================================== %


\newpage


% ********************************************************************* %
%                           ENTRAÎNEMENT                                % 
% ********************************************************************* %
% ================ Métadonnées sur l'entraînement ===================== %

\titreEntrainementFacultatif	{Analyse d'un profil d'énergie potentielle}
\hauteurLargeurCadreReponse		{7mm}{2cm}
\dureeResolutionFacultative		{1} % 1, 2, 3 ou 4
\typeEntrainement				{QCM} % Newton ou QCM ou AN ou vide (exercice "normal")
\nombreColonnesQuestions		{2} % vide, 1, 2, 3, etc.
\avecPlusieursQuestions			{Y} % Y ou N
\initialisationEntrainement
% ===================================================================== %

On reprend le profil d'énergie potentielle de l'entraînement précédent.

Pour chacune des positions suivantes, déterminer si elle est stable ou instable, et si le mouvement au voisinage de ces positions est périodique et/ou harmonique, en indiquant laquelle des réponses suivantes est la bonne. 	\begin{listeQCM2Colonnes}
	\item équilibre stable
	\item équilibre instable
	\item mouvement périodique
	\item mouvement harmonique
\end{listeQCM2Colonnes}

\bigskip

{\itshape Plusieurs bonnes réponses sont possibles.}
\smallskip

% =============================================================================== %
\debutEntrainement
% =============================================================================== %

% ^^^^^^^^^^^^^^^^^^^^^^^^^^^^^^^^^^^^^^ %
%               Question                 %
% ^^^^^^^^^^^^^^^^^^^^^^^^^^^^^^^^^^^^^^ %

\begin{enonce}
	Voisinage de $x^*_1$
\end{enonce}

\reponse{\reponseA{}, \reponseC{} et \reponseD{}}

% ************************************** %



% ^^^^^^^^^^^^^^^^^^^^^^^^^^^^^^^^^^^^^^ %
%               Question                 %
% ^^^^^^^^^^^^^^^^^^^^^^^^^^^^^^^^^^^^^^ %

\begin{enonce}
	Voisinage de $x^*_2$
\end{enonce}

\reponse{\reponseB{}}

% ************************************** %


% ^^^^^^^^^^^^^^^^^^^^^^^^^^^^^^^^^^^^^^ %
%               Question                 %
% ^^^^^^^^^^^^^^^^^^^^^^^^^^^^^^^^^^^^^^ %

\begin{enonce}
	Voisinage de $x^*_3$
\end{enonce}

\reponse{\reponseA{}, \reponseC{} et \reponseD{}}

% ************************************** %


% ^^^^^^^^^^^^^^^^^^^^^^^^^^^^^^^^^^^^^^ %
%               Question                 %
% ^^^^^^^^^^^^^^^^^^^^^^^^^^^^^^^^^^^^^^ %

\begin{enonce}
	Région entre $x_2$ et $x_3$
\end{enonce}

\reponse{\reponseA{} et \reponseC{}}

\begin{corrige}
Le mouvement entre $x_2$ et $x_3$ correspond à un état lié : c'est un mouvement dans un puits de potentiel. Comme le mouvement est à un degré de liberté, il est également périodique. Cependant, les positions extrêmes étant éloignées de la position moyenne (d'équilibre $x^*_3$), ce mouvement n'est pas harmonique.
\end{corrige}

% ************************************** %

% =============================================================================== %
\finEntrainement
% =============================================================================== %




% ********************************************************************* %
%                           ENTRAÎNEMENT                                % 
% ********************************************************************* %
% ================ Métadonnées sur l'entraînement ===================== %

\titreEntrainementFacultatif	{Vitesse à l'infini}
\hauteurLargeurCadreReponse		{8mm}{5cm}
\dureeResolutionFacultative		{1} % 1, 2, 3 ou 4
\typeEntrainement				{AN} % Newton ou QCM ou AN ou vide (exercice "normal")
\nombreColonnesQuestions		{1} % vide, 1, 2, 3, etc.
\avecPlusieursQuestions			{N} % Y ou N
\initialisationEntrainement
% ===================================================================== %

On considère le profil d'énergie potentielle des deux entraînements précédents.\\

Un point matériel de masse $m=\np[kg]{2,30}$ est abandonné avec l'énergie $E_3=\np[kJ]{1,30}$.

% =============================================================================== %
\debutEntrainement
% =============================================================================== %

% ^^^^^^^^^^^^^^^^^^^^^^^^^^^^^^^^^^^^^^ %
%               Question                 %
% ^^^^^^^^^^^^^^^^^^^^^^^^^^^^^^^^^^^^^^ %

\begin{enonce}
Calculer la vitesse du point matériel à l'infini
\end{enonce}

\reponse{$\np[m/s]{33,6}$}

\begin{corrige}
On a vu précédemment que les trajectoires correspondant à l'énergie mécanique $E_3$ sont des états de diffusion, le point matériel peut donc bien s'échapper à l'infini.

Le mouvement du point étant conservatif, on applique la conservation de l'énergie mécanique entre le départ et \glm{l'arrivée} à l'infini : on a
\[
E_3 = \frac{1}{2}mv_\infty^2 
\qquad\text{donc}\qquad
 v_\infty = \sqrt{\frac{2E_3}{m}} = \sqrt{\frac{2\times \SI{1300}{kg.m^{2}.s^{-2}}}{\np[kg]{2,3}}} = \np[m \cdot s^{-1}]{33,6}.
\]
\end{corrige}

% ************************************** %

% =============================================================================== %
\finEntrainement
% =============================================================================== %







% ---------------------------------------------------------------------------- %
%                       Affichage des réponses mélangées                       %
% ---------------------------------------------------------------------------- %
\afficheReponsesMelangees
% ---------------------------------------------------------------------------- %

%%%%%%%%%%%%%%%%%%%%%%%%%%%%%%%%%%%%%%%%%%%%%%%%%%%%%%%%%%%%%%%%%%%%%%%%%%%%%%%%%%%%%%%%%%%%%%
\finFicheEntrainement                                                                        %
%%%%%%%%%%%%%%%%%%%%%%%%%%%%%%%%%%%%%%%%%%%%%%%%%%%%%%%%%%%%%%%%%%%%%%%%%%%%%%%%%%%%%%%%%%%%%%


% ======================================================= % 
% ============== Gestion de la compilation ============== %
% ======================================================= % 
\ifdefined\mainIsLoaded\else
\printReponsesEtCorriges
\end{document}\fi
% ======================================================= % 
% ======================================================= % 